\section{选型指南}
先确定读者需要比较的是数值大小、相对零的变化、组成还是不确定性，再选择条形。
同一比较组应共用范围与宽度，并在标题或表头中写明单位。

\begin{center}
\begin{tabularx}{\linewidth}{@{}>{\centering\arraybackslash}p{23mm}>{\centering\arraybackslash}p{44mm}>{\centering\arraybackslash}X@{}}
\toprule
\multicolumn{1}{c}{\textbf{类型}} & \multicolumn{1}{c}{\textbf{适用情形}} & \multicolumn{1}{c}{\textbf{选择与说明}}\\
\midrule
哑铃图 & 前后值、两组实验对比 & 用 \key{\graddumbbell} 比较位置和差距；两端必须为同一单位。\\
趋势线 & 等间隔观测的变化形状 & 用 \key{\gradspark}；跨行比较绝对波动时共用固定纵向范围。\\

普通条 & 准确率、数量、耗时、完成量 &
用 \key{\gradbar} 表示一个非负值。越长表示越大；耗时越低越好，不能把条长直接理解为优劣。\\
正负条 & 增减幅度、差值、偏差 &
用 \key{min<0} 和正的 \key{max} 显示零点两侧；同列必须共用同一个零点位置。\\
堆叠条 & 互斥类别的构成、资源分配 &
用 \key{\gradstack}；同单位、可相加的数据可分别向零点两侧累计。不同单位或重复计数不应拼成同一组成条。\\
误差线 & 点估计及不确定区间 &
在普通条或正负条上设置 \key{error} 或 \key{error minus}/\key{error plus}。
在表注中明确区间是标准差、标准误还是置信区间。\\
\bottomrule
\end{tabularx}
\end{center}

\subsection{组合使用时的语义}
\key{target} 是一个参考值，\key{band} 是预先给定的参考范围，
它们本身不表示估计误差。误差参数表示从点估计向两侧延伸的非负距离，
不是区间的绝对端点；例如估计 60、区间 $[55,68]$，应设置
\key{error minus=5,error plus=8}。

若希望比较不同方案的总量，保留共同的 \key{max}；若要比较百分比组成，
先由数据处理步骤计算各段百分比，再用统一的 \key{max=100}。
仅设置 \key{segment labels=percent} 会改变标签，不会把原始数据条变成满条。

\subsection{窄列与打印的取舍}
双栏论文优先使用简短数值标签和共享图例；长类别名可放在左侧文本列。
黑白打印同时保留名称、数值或目标位置，避免仅靠颜色表达类别。
幻灯片应增大轨道高度与字体；标签外移后，应为其留出额外垂直空间。
